% !PRINTED 636
\weberpage{636}
%<<CONT
erheben und dann alles multipliciren, und wenn wir noch beachten, dass
$\varepsilon^{g^{n-1}} = \varepsilon$ ist, so heben sich im Product der
linken Seite der Gleichungen (5) alle Resolventen mit Ausnahme von
$(\varepsilon, x)$ heraus, und es folgt

\weq{8}{(\varepsilon, x)^{1-g^{n-1}} = f_0^{g^{n-2}} f_1^{g^{n-3}} \cdots f_{n-2}.}

Bezeichnen wir nun mit $\lambda$ einen noch unbestimmten ganzzahligen
Exponenten, so leiten wir aus (8) die folgende Relation her:

\weq{9}{(\varepsilon^\lambda, x)^n = \left[ (\varepsilon, x)^{\frac{g^{n-1}-1}{n}} (\varepsilon^\lambda, x) \right]^n f_0^{g^{n-2}} f_1^{g^{n-3}} \cdots f_{n-2}.}

Verstehen wir unter $g_1$ irgend eine feste primitive Wurzel von $n$, so
können wir $g = g_1 + l n$ setzen, worin $l$ eine beliebige ganze Zahl
bedeutet, da ja die primitiven Wurzeln nur bis auf Vielfache von $n$
definirt sind. Dann ist

\[
g^{n-1} \equiv g_1^{n-1} + n(n-1) l g_1^{n-2} \pmod{n^2},
\]

woraus

\[
\frac{g^{n-1}-1}{n} \equiv \frac{g_1^{n-1}-1}{n} - l g^{n-2} \pmod{n},
\]

und $l$ lässt sich so bestimmen, dass

\[
\frac{g^{n-1}-1}{n} \equiv 1 \pmod{n}
\]

wird.

In (9) setzen wir nun unter dieser Voraussetzung

\weq{10}{\lambda = \frac{g^{n-1}-1}{n} \equiv 1 \pmod{n},}

und wie in (4)

\weq{11}{F(\varepsilon) = (\varepsilon^\lambda, x)(\varepsilon, x)^{-\lambda} = (\varepsilon, x)^{nk},}

worin $\varepsilon^\lambda = \varepsilon$ und $1 - \lambda$, was durch $n$
theilbar ist, $= nk$ gesetzt ist. Daraus erhalten wir

\weq{12}{(\varepsilon, x)^n = [F(\varepsilon)]^n f_0^{g^{n-2}} f_1^{g^{n-3}} \cdots f_{n-2}.}

Diese Formeln lassen sich verallgemeinern, wenn man beachtet, dass die
Functionen $f_0, f_1 \ldots f_{n-2}$ eine cyklische Permutation erleiden,
wenn die Substitution $(\varepsilon, \varepsilon^g)$ ausgeführt wird.
Setzen wir also fest, dass $f_h = f_k$ sein soll, wenn $h \equiv k
\pmod{n-1}$ ist, und setzen

\weq{13}{F(\varepsilon^{g^\nu}) = F_\nu = (\varepsilon^{g^\nu}, x)^{nk},}

so ergiebt die Substitution $(\varepsilon, \varepsilon^{g^\nu})$ in (12)

\weq{14}{(\varepsilon^{g^\nu}, x)^n = F_\nu^n f_\nu^{g^{n-2}} f_{\nu+1}^{g^{n-3}} \cdots f_{\nu+n-2}.}
